С развитием цифровых технологий и повсеместным распространением мобильных устройств
взаимодействие горожан с городской средой приобретает новые формы.
Жители, ежедневно перемещающиеся по городу, первыми замечают изменения инфраструктуры,
аварийные ситуации и иные события, требующие оперативного реагирования со стороны
муниципальных служб. В то же время оперативная фиксация подобных событий и передача
информации ответственным органам остаётся затруднённой: гражданину необходимо
самостоятельно определить ответственную службу, найти её контактные данные и составить
обращение по телефону или электронной почте. Единый цифровой канал для подачи обращений,
как правило, отсутствует, а обратная связь о статусе обращения не предоставляется.

Всё это делает актуальной разработку программной системы фиксации
городских событий. Такая система позволит жителям быстро сообщать о проблемах
городской инфраструктуры, прикладывая геолокацию, фотографии и
текстовое описание, а диспетчерам~--- оперативно обрабатывать поступившие обращения
с помощью интерактивной карты и автоматической маршрутизации к ответственным службам.

Актуальность подтверждается следующими факторами:
\begin{itemize}
    \item растущие требования к оперативности сбора и обработки обращений граждан
          по вопросам городской инфраструктуры;
    \item необходимость обеспечения прозрачности процесса обработки обращений и
          информирования заявителей о текущем статусе;
    \item потребность муниципальных служб в централизованном приёме и диспетчеризации
          обращений, поступающих от граждан;
    \item повсеместное распространение мобильных устройств, позволяющих фиксировать
          проблемы непосредственно на месте их обнаружения;
    \item отсутствие единого решения, объединяющего фиксацию проблемы, диспетчеризацию
          на карте и маршрутизацию к ответственной службе, применимого к городу Кирову.
\end{itemize}

Целью преддипломной практики является обобщение результатов проектирования и разработки
программной системы фиксации городских событий, выполненной в рамках двух курсовых работ,
а также подготовка материалов для выпускной квалификационной работы.

Программная система включает три основных компонента:
\begin{enumerate}
    \item серверную часть~--- программный интерфейс на платформе ASP.NET Core с базой
          данных PostgreSQL и расширением PostGIS для геопространственных запросов;
    \item диспетчерскую веб-панель~--- автоматизированное рабочее место оператора
          на React с интерактивной картой обращений;
    \item кроссплатформенное мобильное приложение для жителей~--- приложение на
          React Native для фиксации проблем городской инфраструктуры.
\end{enumerate}

Для достижения данной цели в ходе практики решались следующие задачи:
\begin{enumerate}
    \item систематизация результатов анализа предметной области и обзора аналогов,
          выполненных на первом и втором этапах разработки;
    \item обобщение проектных решений по структуре системы, структуре базы данных,
          пользовательскому интерфейсу веб-панели диспетчера и мобильного приложения;
    \item описание реализации серверной части, веб-панели диспетчера и мобильного
          приложения с обоснованием выбора программных средств;
    \item обобщение результатов тестирования всех компонентов программной системы;
    \item подготовка текста пояснительной записки к выпускной квалификационной работе.
\end{enumerate}
